\patentsectionstart{权 利 要 求 书}{1}

\begin{claimitem}
\item 一种人机共享记忆的自愈索引与按需下探方法，其特征在于，包括：
\begin{enumerate}[label=(\arabic*),leftmargin=2.6em,itemsep=0.65em,topsep=0.55em]
\item 在文件系统中保存结构化地图文件以及与记忆对象对应的对象资料包目录，所述对象资料包目录包括正文文件，所述结构化地图文件和对象资料包目录构成人类处理端与人工智能处理端共用的文件事实源；
\item 在建立卡片索引或正文文件被保存后，为所述正文文件生成或刷新卡片条目，所述卡片条目包括路径、标题、摘要和存储指纹，且不保存正文全文；所述存储指纹包括作为内容版本标记的正文内容摘要值、修改时间和文件大小；
\item 在接收到上下文请求时，将正文文件的当前修改时间和当前文件大小分别与所述卡片条目中保存的修改时间和文件大小比较；当二者均相等时，复用所述卡片条目的摘要；当至少一者不相等或对应卡片条目不存在时，仅读取对应正文文件并刷新所述卡片条目；
\item 依据当前节点、路线状态和文本相关性对记忆对象排序，在预设上限内输出结构化地图投影、命中卡片摘要、目标路径以及对应的召回理由；在接收到针对所述目标路径的显式读取请求时，读取对应的正文文件。
\end{enumerate}

\item 根据权利要求 1 所述的方法，其特征在于，还包括：接收对所述结构化地图文件的写入命令，校验所述写入命令的命令标识和基线版本号；当所述基线版本号落后于当前版本号时，比较所述写入命令的字段触达范围与基线版本之后已提交命令的字段触达范围，在二者不重叠时基于当前版本应用所述写入命令，在二者重叠时返回冲突路径、当前已提交数据和待写入数据；对通过校验的写入命令予以持久化。

\item 根据权利要求 1 所述的方法，其特征在于，经由本系统提供的正文保存接口保存正文文件时，保存请求携带调用方先前读取到的基线内容版本标记；在目标路径的写入锁内，将所述基线内容版本标记与当前正文文件的内容版本标记比较；二者不一致时，不覆盖当前正文文件并返回当前正文文件及其内容版本标记；二者一致时，通过临时文件原子替换所述正文文件并刷新对应卡片条目。

\item 根据权利要求 1 所述的方法，其特征在于，由人类经由本系统提供的正文保存接口保存正文文件时，向追加式日志写入未确认记录，所述未确认记录包括路径、内容版本标记、修改时间、正文片段和记录时间；按路径重放所述追加式日志并以同一路径的最后一条状态记录确定该路径是否未确认；将结构标注和未确认记录合并为人类新输入列表，在确认接口返回确认结果前持续在会话上下文中呈现所述未确认记录；同一路径再次由人类经由所述正文保存接口保存时，写入新的未确认记录。

\item 根据权利要求 1 所述的方法，其特征在于，所述对象资料包目录还包括资产文件，所述卡片条目还包括附件清单；所述资产文件包括图片、文档或多媒体文件，并在上下文中以路径和元数据表示，二进制内容仅在针对资产路径的显式受控读取请求中返回，且由具备相应读取能力的调用方处理。

\item 根据权利要求 1 所述的方法，其特征在于，所述排序包括：将当前节点及其图邻域作为优先对象，为未解决的问题节点和待验证方案增加状态优先级，并对卡片摘要执行 BM25 或其他基于词频统计的文本相关性评分；所述召回理由至少包括命中当前节点、图邻域关系、状态优先级或文本相关性评分中的一种。

\item 根据权利要求 2 所述的方法，其特征在于，对所述结构化地图文件的持久化包括：根据命令标识和依据命令内容计算的请求摘要判断所述写入命令是否已经提交，对命令标识和请求摘要均相同的重复写入返回已提交结果，对命令标识相同而请求摘要不同的写入予以拒绝；应用所述写入命令并校验所得地图；将包含命令标识、请求摘要、原版本校验值、字段触达范围、目标版本、所得地图校验值和所得地图的准备记录写入预写日志；通过临时文件原子替换所述结构化地图文件；写入提交记录并返回新版本号和所得地图校验值；系统启动时，对具有准备记录而没有提交记录的写入，若磁盘地图为所述准备记录对应的原版本，则写入准备记录中的所得地图并补写提交记录，若磁盘地图已为准备记录中的所得地图，则补写提交记录，若磁盘地图既不等于所述原版本也不等于所述所得地图，则保存隔离副本并返回恢复冲突。

\item 一种人机共享记忆的自愈索引与按需下探系统，其特征在于，包括：
\begin{enumerate}[label=(\arabic*),leftmargin=2.6em,itemsep=0.65em,topsep=0.55em]
\item 文件存储模块，用于保存结构化地图文件和对象资料包目录，并使所述结构化地图文件和对象资料包目录作为人类处理端与人工智能处理端共用的文件事实源；
\item 地图命令模块，用于校验地图写入命令的命令标识和基线版本号，在基线版本落后时根据字段触达范围判断是基于当前版本重放还是返回重叠冲突，并通过预写日志、原子替换和提交记录持久化通过校验的地图写入命令；
\item 卡片索引模块，用于为正文文件生成或刷新包括路径、标题、摘要和存储指纹且不含正文全文的卡片条目，所述存储指纹包括作为内容版本标记的正文内容摘要值、修改时间和文件大小；
\item 新鲜度校验模块，用于比较正文文件与卡片条目中保存的修改时间和文件大小，在二者均相等时复用摘要，在至少一者不相等或卡片条目不存在时仅读取对应正文文件并刷新卡片条目；
\item 上下文投影模块，用于按照当前节点、路线状态和文本相关性在预设上限内输出结构化地图投影、卡片摘要、目标路径和召回理由；
\item 下探读取模块，用于响应针对目标路径的显式读取请求而读取正文文件。
\end{enumerate}

\item 一种电子设备，包括存储器和处理器，所述存储器中存储有计算机程序，其特征在于，所述处理器执行所述计算机程序时实现权利要求 1 至 7 中任一项所述的方法。

\item 一种计算机可读存储介质，其上存储有计算机程序，其特征在于，所述计算机程序被处理器执行时实现权利要求 1 至 7 中任一项所述的方法。
\end{claimitem}
